\section{Introduction}
\label{sec:intro}

Electronic health records (EHRs) generate administrative data on every
patient encounter: International Classification of Diseases (ICD)
diagnosis codes, procedure codes, and medication orders. A large
secondary-use literature reuses this data to answer the question
\emph{does a patient have clinical condition $Z$?} This task is
\emph{electronic phenotyping}: inferring a latent clinical state from
coarse administrative signals \citep{banda2018phenotyping,
hripcsak2013correlating}. Phenome-wide association studies
\citep{denny2010phewas}, the Electronic Medical Records and Genomics
network \citep{gottesman2013emerge,newton2013validation}, the PheKB
algorithm catalog \citep{kirby2016phekb}, anchor-and-learn classifiers
\citep{halpern2016anchor}, high-throughput phenotyping
\citep{liao2015highthroughput}, semi-supervised methods such as
PheNorm \citep{yu2018phenorm} and PheCAP \citep{zhang2019phecap}, and
probabilistic phenotypes \citep{pivovarov2015learning} all build
phenotype definitions on top of diagnosis codes and related EHR
features.

The difficulty is that the diagnosis code is a noisy and
\emph{systematically biased} proxy for the true clinical state. A code
is generated only when a clinician records an encounter, when a billing
process attaches a code to that encounter, and when coverage and
reimbursement incentives make that code worth recording. The fraction
of patients carrying a code for condition $Z$ is therefore not the
prevalence of $Z$. Code sensitivity (the probability that a true case
is coded) and code specificity (the probability that a non-case is not
coded) are both unknown, and both depend on factors, encounter
frequency, disease severity, and local coding practice, that have
nothing to do with the underlying epidemiology
\citep{hripcsak2013correlating,carrell2014surveillance}.

\subsection{The bridge to masked-data inference}

The central observation of this paper is that electronic phenotyping
is mathematically isomorphic to the masked-data series-system
identifiability problem of \citet{towell2026masked}. In that
framework, a latent cause is masked, but a \emph{candidate set}
containing it is observed; coarsening-at-random conditions
\citep{heitjan1991ignorability,gill1997coarsening} on the masking
mechanism determine whether the cause-level parameters can be
recovered.

The translation to phenotyping is direct. The true clinical state
$\Y \in \{0,1\}$ is the latent cause. The set of clinical states
consistent with a patient's observed diagnosis codes is the candidate
set. The coding process, $\Prob(\text{code} \mid \Y)$, shaped by
clinician behavior and billing, is the masking mechanism. A
chart-reviewed (clinician-adjudicated) patient, for whom the true
state is established by manual record review, is a singleton candidate
set: the truth is known exactly. The C1-C2-C3 conditions for ignorable
coarsening port directly. This framework has been applied to scRNA-seq
zero inflation \citep{towell2026scrnacoarsening} and to spatial
transcriptomics deconvolution \citep{towell2026spatialcoarsening};
electronic phenotyping is the present application.

The phenotyping case has a distinctive feature. In the transcriptomics
applications the coarsening mechanism is approximately non-informative
once technical covariates are modeled. In phenotyping it is
\emph{structurally} informative. Sicker patients have more encounters
and are coded more; billing incentives change which codes are recorded;
surveillance and detection intensity track disease severity
\citep{haut2011surveillance}. The coarsening-at-random condition that
fails is C2, non-informative coding. This single observation explains
why code-only prevalence estimates are biased and why chart review is
a necessity rather than a luxury.

\subsection{Prior work on latent-class phenotyping}

Estimating disease status without a gold standard has a deep prior
literature that we do not claim to originate.
\citet{rogan1978estimating} gave the canonical
sensitivity-and-specificity correction for apparent prevalence, the
plug-in estimator that recovers $\pi$ from the observed code
frequency once $\mathrm{sens}$ and $\mathrm{spec}$ are known; this
paper's chart-review estimator is that estimator with the two error
rates supplied by the adjudicated subsample. \citet{hui1980estimating}
introduced the latent-class model for estimating diagnostic-test error
rates when no perfect reference test exists, and
\citet{dawid1979maximum} gave the EM algorithm for rater-error rates
that is the methodological ancestor of all such work. The
verification-bias literature (\citealp{begg1983assessment}; see
\citealp{hubbard2020outcome} for a recent phenotyping-specific
treatment) recognized that when verification of disease status is
selective, naive estimates of test accuracy and prevalence are biased;
the case-mix gap of \cref{cor:casemix} is the coarsening recasting of
this phenomenon. eMERGE
\citep{gottesman2013emerge} and PheKB \citep{kirby2016phekb} are
large-scale phenotyping efforts, and the anchor-and-learn framework
\citep{halpern2016anchor} uses high-precision signals in a way
structurally analogous to singleton candidate sets. Our contribution
is not latent-class phenotyping itself. It is the masked-cause
unification: a single coarsening vocabulary that classifies the
failure modes of code-based phenotyping through the C1-C2-C3
conditions, pinpoints informative coding as the C2 violation, and
quantifies the resulting bias.

\subsection{Contributions}

We contribute:

\begin{itemize}
\item \textbf{A bridge from electronic phenotyping to masked-data
  inference} (\cref{sec:translation}), with an explicit translation
  table mapping the coding mechanism to the masking mechanism and a
  chart-reviewed patient to a singleton candidate set.

\item \textbf{A glass-ceiling theorem} (\cref{thm:glass-ceiling},
  \cref{sec:identifiability}): with an unrestricted coding mechanism,
  prevalence, code sensitivity, and code specificity are jointly
  non-identifiable from code data alone. This formalizes why the
  fraction of patients carrying a code is a biased prevalence
  estimate.

\item \textbf{An identifiability theorem with chart review}
  (\cref{thm:identifiability-chart}, \cref{sec:identifiability}): a
  chart-reviewed subsample covering both diseased and non-diseased
  patients identifies the coding model, and cohort prevalence then
  becomes identifiable. This is the direct analog of the spike-in
  identifiability theorem of \citet{towell2026scrnacoarsening}.

\item \textbf{A code-frequency consistency theorem}
  (\cref{thm:code-total}, \cref{sec:identifiability}): a fitted
  phenotype model reproduces the marginal observed code frequencies
  exactly at an interior MLE. Marginal-fit checks therefore cannot
  detect prevalence bias.

\item \textbf{A bias bound under informative coding}
  (\cref{thm:bias-informative}, \cref{sec:methodology}): when code
  sensitivity correlates with latent disease severity, a C2 violation,
  the code-only prevalence estimator has bias bounded by a function of
  the severity-coding correlation. The chart-review-calibrated
  estimator's residual bias is bounded by the case-mix gap between the
  chart-reviewed subsample and the cohort, the analog of the
  ERCC-versus-endogenous capture-efficiency gap in
  \citet{towell2026scrnacoarsening}.

\item \textbf{Simulation validation} (\cref{sec:validation}), with a
  MIMIC-IV \citep{johnson2023mimiciv} real-data application described.
\end{itemize}

The paper's central message: code-based prevalence estimation is
biased by structure, not by accident, because coding is informative
and C2 fails. The bias is not directional in a simple way: depending
on the balance between code sensitivity and false-positive inflation,
code-only prevalence can understate the truth, overstate it, or pass
through zero as the severity-coding correlation varies (see
\cref{sec:validation}, \cref{tab:informative}). Chart review supplies
the singleton candidate sets that restore identifiability, and its
value is quantified by how well the reviewed subsample represents the
cohort.
